% =====================================================================
% Paper skeleton (cloud task T4, docs/cloud_tasks/T4_paper_skeleton.md)
% IEEE RA-L style, IEEEtran journal class.
%
% SOURCE RULES (T4): no invented numbers or references. Every number
% carries a \src{} tag naming the file it comes from (key list below);
% anything without a source is a \TODO{}.
%
% Inputs that were NOT available when this skeleton was written:
%   - docs/cloud_tasks/handover.md            (absent -> every handover item is \TODO)
%   - docs/cloud_tasks/T1_related_work_report.md (absent -> Related Work is a \TODO list)
%   - the results report of the frozen p9d_analyze.py (absent; only the
%     post-check report is in eval/phase9/P9-D/)
%   - the P9-E specification of T1, T2, F1--F4 (absent; placeholders only)
%
% Source keys used by \src{}:
%   T4  docs/cloud_tasks/T4_paper_skeleton.md   (claims list, provided by Marco)
%   CR  eval/phase9/P9-D/criterion.md            (pre-registered criterion)
%   A1  eval/phase9/P9-D/addendum_A1.md
%   A2  eval/phase9/P9-D/addendum_A2.md
%   PR  eval/phase9/P9-D/prerun.md
%   PCR eval/phase9/P9-D/postcheck_report.md
%   PCT eval/phase9/P9-D/p9d_postcheck.md        (post-check task, frozen verdict table)
%   CX  eval/phase9/P9-D/p9d_crowd_explore.md    (exploratory task spec; no results)
% =====================================================================
\documentclass[journal]{IEEEtran}

\usepackage[utf8]{inputenc}
\usepackage[T1]{fontenc}
\usepackage{amsmath,amssymb}
\usepackage{graphicx}
\usepackage{booktabs}
\usepackage{adjustbox}
\usepackage{multirow}
\usepackage{xcolor}
\usepackage{siunitx}
\usepackage{url}
\usepackage{cite}

% ---- draft macros ----------------------------------------------------
\newcommand{\TODO}[1]{{\color{red}\textbf{[TODO:} #1\textbf{]}}}
\newcommand{\src}[1]{\textsuperscript{\color{blue}\scriptsize[#1]}}
\newcommand{\figplaceholder}[2]{%
  \fbox{\parbox[c][#2][c]{0.95\linewidth}{\centering\small\color{gray}#1}}}

\begin{document}

\title{\TODO{Title}: Pedestrians Inside Visual-Inertial SLAM ---
Online Estimation Without Measurable Camera Cost, and Why They Do Not Help Odometry}

\author{\TODO{Authors and affiliations (handover)}%
\thanks{\TODO{Manuscript dates, funding (handover).}}}

\markboth{IEEE Robotics and Automation Letters. Preprint version. \TODO{date}}%
{\TODO{Short author list}: Pedestrians Inside Visual-Inertial SLAM}

\maketitle

% =====================================================================
\begin{abstract}
We place pedestrians inside the local bundle adjustment of a mono-inertial
visual SLAM system (an ORB-SLAM3-based fork) as persistent world-frame states,
estimated online alongside the camera.
In a pre-registered evaluation on \TODO{dataset name and sequences (handover)},
the pedestrian layer estimates pedestrian velocity better than a per-track
Kalman filter on the same measurements, and its position is no worse than
per-frame back-projection; this result rests on the best-aligned runs, and a
post-hoc sensitivity analysis is pending\src{T4}.
On the camera we find no evidence of degradation: the ratio of median ATE with
and without the layer is 1.00 [0.94,\,1.06] and 1.07 [0.95,\,1.18] on the two
confirmatory sequences\src{T4}; the pre-registered non-inferiority test passes
at a 10\% margin on one sequence and is inconclusive on the other\src{T4,\,A2}.
We then explain why pedestrians cannot improve odometry under normal
conditions: camera-relative pedestrian factors are absorbed exactly, and in
the regime where every keyframe sees at least two static points the
pedestrians carry no camera information. Threaded throughput is about
6--7$\times$ real time\src{PCR\,\S5}.
\TODO{Final abstract once results report and handover are available.}
\end{abstract}

\begin{IEEEkeywords}
\TODO{RA-L keywords, e.g.\ SLAM; visual-inertial odometry; human detection and tracking}
\end{IEEEkeywords}

% =====================================================================
\section{Introduction}
\label{sec:intro}

\IEEEPARstart{M}{otivation}: \TODO{(handover) Why estimate pedestrians
inside the SLAM back-end; the common expectation that dynamic agents
could also constrain the camera; what is missing in prior work (from T1 report,
absent).}

\TODO{(handover) One paragraph on the system: pedestrian states in the local
BA of the ORB-SLAM3-based \texttt{visual\_sgraphs} fork; deterministic mode;
evaluation design (pre-registered criterion, two independent hypothesis
families\src{A2\,\S A2.3}).}

\paragraph*{Contributions}
The claims below are stated exactly at the strength of the evidence; none is
upgraded beyond its status in the claims list\src{T4}.
\begin{enumerate}
  \item \emph{System.} Pedestrians as persistent states inside the local bundle
  adjustment of a mono-inertial visual SLAM system, estimated online.
  The threaded configuration runs at about 6--7$\times$ real time in
  throughput (wall-clock \SIrange{99.8}{121.6}{\second} for
  \SI{738}{\second} of sensor data, one run per arm)\src{PCR\,\S5};
  latency was not measured.
  \item \emph{Pedestrian estimation (pre-registered, supported).}
  H2a: pedestrian velocity is better than a per-track Kalman filter on the same
  measurements; H2b: pedestrian position is no worse than per-frame
  back-projection\src{T4,\,PCT}. \emph{Caveat:} H2 rests on the best-aligned
  runs (the pre-registered map-inclusion filter keeps 6--12 of 60 subway runs
  per cell and arm)\src{PCR\,\S1}; a post-hoc sensitivity analysis is
  pending\src{T4}.
  \item \emph{Camera.} No evidence of degradation: median-ATE ratio
  1.00 [0.94,\,1.06] and 1.07 [0.95,\,1.18]\src{T4}; non-inferior at a 10\%
  margin on one sequence, inconclusive on the other; the pre-registered
  camera hypothesis H1 formally failed\src{T4,\,PCT}.
  \item \emph{Analysis: why pedestrians do not help the camera.}
  (i) exact absorption of camera-relative pedestrian factors;
  (ii) an information analysis showing that with $\geq 2$ static points per
  node the pedestrians bring no camera benefit;
  (iii) unidentifiability of stature, resolved by MAP shrinkage;
  (iv) an erratum on range: depth-$z$ versus ray length;
  (v) the chaotic run-to-run spread of ORB-SLAM3 and its effect on the effect
  sizes that can be detected\src{T4}.
\end{enumerate}

\paragraph*{What we do not claim}
We do not claim a camera improvement; we do not claim floor levels (the
pre-registered levels hypothesis H4 failed\src{PCT}; levels are left to a
follow-up paper\src{T4}); and we do not claim pedestrian accuracy on real
data\src{T4}.

% =====================================================================
\section{Related Work}
\label{sec:related}

\TODO{The T1 related-work report (\texttt{docs/cloud\_tasks/T1\_related\_work\_report.md})
is absent. No related work is written and no reference is cited until it exists.
Structure to fill from T1, one paragraph per group, each method positioned
against our system (online or offline; sensors; humans inside the back-end
or not; human representation; camera-improvement claim; pedestrian metrics;
datasets; code).}

\begin{itemize}
  \item \TODO{Baselines named in T4 itself: ORB-SLAM3 (the base system);
  per-track Kalman filter (B1); per-frame back-projection (B0).}
  \item \TODO{3D dynamic scene graphs with humans: Kimera / 3D Dynamic Scene
  Graphs (uHumans datasets), Hydra.}
  \item \TODO{Human mesh / trajectory recovery with camera motion: BodySLAM,
  SLAHMR, PACE, WHAM, TRAM, GLAMR, JOSH.}
  \item \TODO{Dynamic-object SLAM: AirDOS, DynoSAM, VDO-SLAM, ClusterVO,
  DynaSLAM.}
  \item \TODO{2023--2026 work estimating humans jointly with a SLAM/VIO back-end
  (from T1 search).}
  \item \TODO{Papers reporting the opposite camera finding (a camera benefit from
  humans or dynamic objects) and their conditions (T1, item 4); contrast with
  Sec.~\ref{sec:analysis}.}
  \item \TODO{Positioning paragraph: which claims are new, incremental, or
  already shown (T1, item 3).}
\end{itemize}
% Method names above are from the T1 task file (docs/cloud_tasks/T1_related_work.md);
% no bib entries are created for them until the T1 report supplies verified sources.

% =====================================================================
\section{Method}
\label{sec:method}

The handover (\texttt{docs/cloud\_tasks/handover.md}) is not available; every
item below that depends on it is marked \TODO{}. Statements without a
\TODO{} are taken from the P9-D reports and carry their source.

\subsection{System overview}
\TODO{(handover) Block diagram and description.} The pedestrian layer lives
inside the local bundle adjustment of ORB-SLAM3 in the \texttt{visual\_sgraphs}
fork, mono-inertial; pedestrian instance masks are applied to the camera
front end in every configuration, including the configuration without the
layer\src{CR\,\S2}. The layer reads per-frame perception rows: person (P) rows,
\TODO{C rows: meaning (handover)}, and optional foot (F) rows\src{A1\,\S A1.2}.
Detections and keypoints come from RTMW\src{PR\,\S1}; tracks from
\TODO{P1-A tracker (handover)}.

\subsection{States}
\TODO{(handover) Per-person world-frame root state: position, velocity;
persistence and creation/retirement of tracks; anchoring to keyframes;
per-track stature parameter $b$ (``per-track b'' is reported
descriptively\src{CR\,\S6}).}
\begin{equation}
  \mathbf{x}_{j}(t) = \TODO{\text{state definition}}
\end{equation}

\subsection{Measurement factors}
\TODO{(handover) Image factors of the person observation: bearing of the root
keypoint, apparent body size (nose$\to$heel) as range cue; noise models,
robust kernels, gating.}
\begin{equation}
  \mathbf{r}_{ij} = \TODO{\text{residual}}
\end{equation}

\subsection{Range from apparent size: depth-$z$, not ray length}
\TODO{(handover) Range model at a fixed stature; the erratum: the size cue
constrains the depth $z$ along the optical axis, not the length of the ray;
effect of the earlier ray-length form and where it was corrected.}

\subsection{Stature}
\TODO{(handover) Stature is not identifiable from monocular size alone (it
trades off against range); a prior on stature gives MAP shrinkage towards the
prior mean; prior values. Cross-reference Sec.~\ref{sec:an-stature}.}

\subsection{Motion prior}
\TODO{(handover) Constant-velocity / white-noise-acceleration prior between
consecutive pedestrian states; process-noise density; behaviour across gaps.}
\begin{equation}
  \mathbf{r}^{\text{mp}}_{j,k} = \TODO{\text{motion-prior residual}}
\end{equation}

\subsection{Footprints}
\TODO{(handover) Stance-foot footprint factors.} The foot contact used by the
contact factors is the midpoint of the left and right foot's lowest valid
keypoint (F rows)\src{PR\,\S1}. F-row coverage on the evaluation sequences is
46.8\%, 54.1\% and 71.2\% of person rows (subway\_s1\_36h, subway\_s1\_24h,
office\_s1\_06h)\src{PR\,\S1}. Configuration \texttt{FloorContact: none} is the
base used by the ON arm\src{CR\,\S2}. \TODO{(handover) Whether footprints are
active in the ON arm; factor form.}

\subsection{Integration into the local BA}
\TODO{(handover) Which keyframes and pedestrian states enter the local window;
fixed vs.\ free variables; marginalisation / dropping of old pedestrian states;
interaction with IMU factors; solver settings; code locations
(\texttt{PedestrianBA.cc}\src{A1\,\S A1.2}).}

\subsection{Determinism}
The evaluation runs ORB-SLAM3 in a deterministic mode, which is not stock
ORB-SLAM3; three uninitialised-read contexts of the vanilla code
remain\src{CR\,\S9}. With the levels configuration, five identical invocations
on subway\_s1\_24h were byte-identical\src{PR\,\S3}. Run-to-run variation is
introduced on purpose by a seeded jitter and a staging path of $40+k$
characters for seed $k$, identical across arms\src{CR\,\S2}; output still
depends on the settings-file text and path length (layout
sensitivity)\src{PCR\,\S4}. \TODO{(handover) What the deterministic mode
changes (thread scheduling, etc.).}

\subsection{Floor levels (not claimed)}
An optional floor-level landmark module (\texttt{FloorContact: levels}, no
platform prior) exists and was evaluated\src{CR\,\S2}; its hypothesis H4
failed\src{PCT} and the module is left to a follow-up paper\src{T4}.
\TODO{Decide whether this subsection stays (one paragraph) or moves to
Limitations.}

% =====================================================================
\section{Analysis: Why Pedestrians Do Not Help the Camera}
\label{sec:analysis}

\subsection{Absorption of camera-relative factors}
\label{sec:an-absorb}
\TODO{(handover) Exact statement and proof: a factor that depends only on the
pedestrian state relative to the observing camera is absorbed by the free
pedestrian state, leaving the camera marginal unchanged.}

\subsection{Information regime}
\label{sec:an-info}
\TODO{(handover) Fisher-information analysis; with $\geq 2$ static points per
node\src{T4} the pedestrians add no camera information; pedestrians carry
camera information only near zero static support.} The exploratory
crowd-stratified analysis that probes the high-occupancy regime on the
existing runs was specified\src{CX} \TODO{its report
(\texttt{crowd\_explore\_report.md}) is absent; include only if it exists, labelled
exploratory and post hoc}.

\subsection{Stature unidentifiability}
\label{sec:an-stature}
\TODO{(handover) Derivation; MAP shrinkage of per-track $b$.}

\subsection{Range erratum: depth-$z$ vs.\ ray length}
\label{sec:an-range}
\TODO{(handover) Size of the error as a function of off-axis angle; which
results it affected.}

\subsection{Chaotic ORB-SLAM3 spread and detectable effect sizes}
\label{sec:an-spread}
Stage~0 gave a relative spread $\sigma_r/\text{median}$ of the SE(3) ATE of
0.227 (subway\_s1\_36h) and 0.198 (subway\_s1\_24h)\src{CR\,\S3}. For a
one-sided $\alpha=0.05$, power 0.8 and a 5\% margin, the required number of
seeds per arm is
\begin{equation}
  N = 2\,(1.645+0.84)^2\,\bigl(\sigma_r/\text{median}\,/\,0.05\bigr)^2,
\end{equation}
i.e.\ 260 and 200 seeds\src{CR\,\S3}. On 20 identical-binary runs entered
twice (as OFF and as a fake ON), the 95\% interval of the median ratio was
already [0.828,\,1.203] and [0.816,\,1.210]\src{PR\,\S4}. With $N=60$ per arm
the detectable margin is about 10\% on M1 and $+0.20$ on the map-event
rate\src{A2\,\S A2.2}. \TODO{(handover) Source of the chaos (thread timing,
layout sensitivity\src{PCR\,\S4}); minimum detectable camera effect vs.\ the
upper bound of any plausible pedestrian benefit from
Sec.~\ref{sec:an-info}.}

% =====================================================================
\section{Experiments}
\label{sec:exp}

\subsection{Data}
\TODO{(handover) uHumans2 description\,/\,citation (from T1).}
Confirmatory sequences: subway\_s1\_36h and subway\_s1\_24h, mono-inertial;
office\_s1\_06h is descriptive only\src{CR\,\S3}. Tuning used only
apartment\_s1\_01h, apartment\_s1\_02h and office\_s1\_12h\src{CR\,\S9}.

\subsection{Protocol (pre-registered)}
Three arms, all deterministic and mono-inertial with masks: OFF (layer
disabled), ON (pedestrian layer, \texttt{FloorContact: none}), LEV (floor
levels)\src{CR\,\S2}. $N=60$ seeds per arm on each subway sequence and 20 on
office, the same seeds for every arm, about 420 runs\src{A2\,\S A2.1,\,A2.5},
interleaved by seed\src{PR\,\S5}.

\emph{Camera metrics:} M1, SE(3) ATE of the main map (a run without a
main-map metric counts as $+\infty$); M2, the map-event indicator ($\geq 1$
reset or new map); secondary: Sim(3) ATE, $|\log\hat s|$, coverage, scale
failure, exit code\src{CR\,\S4}.
\emph{Pedestrian metrics:} per run, the median camera-frame position error and
the median world-frame velocity error of the graph, B0 (per-frame
back-projection), B1-causal (per-track Kalman filter) and B1-RTS, on the same
(keyframe, track) samples\src{CR\,\S4,\,PR\,\S4}.
\emph{Map inclusion} for pedestrian samples: $|\log\hat s|<0.2$, Sim(3)
camera-centre RMSE $\leq\SI{0.5}{\metre}$, extent $\geq\SI{2}{\metre}$,
$\geq 50$ keyframes\src{CR\,\S4}.

\emph{Hypotheses} (one-sided $\alpha=0.05$; two independent families, each a
fixed sequence\src{A2\,\S A2.3}):
H1, ON vs.\ OFF, upper bound of the M1 median ratio $<1.10$ and of the M2
difference $<+0.20$\src{A2\,\S A2.2};
H2a, Wilcoxon signed-rank of the graph/B1-causal velocity log-ratio $<0$;
H2b, upper bound of the median graph/B0 position ratio $<1.05$\src{CR\,\S5,\,A2\,\S A2.2};
H3, LEV vs.\ OFF as H1; H4, LEV vs.\ ON position\src{CR\,\S5}.

\subsection{Camera results}
\begin{table}[t]
\centering
\caption{\TODO{P9-E T1 (camera).} Per-cell point estimate and two-sided 95\%
interval for each camera quantity\src{A2\,\S A2.4}.
\TODO{Exact layout per P9-E spec (absent).}}
\label{tab:camera}
\begin{adjustbox}{max width=\columnwidth}
\setlength{\tabcolsep}{3pt}\footnotesize
\begin{tabular}{@{}llccc@{}}
\toprule
Cell & Quantity & ON/OFF & LEV/OFF & Verdict \\
\midrule
\multirow{2}{*}{subway\_s1\_36h} & M1 ratio [95\%] & \TODO{} & \TODO{} & \TODO{} \\
                                 & M2 diff.\ [95\%] & \TODO{} & \TODO{} & \TODO{} \\
\multirow{2}{*}{subway\_s1\_24h} & M1 ratio [95\%] & \TODO{} & \TODO{} & \TODO{} \\
                                 & M2 diff.\ [95\%] & \TODO{} & \TODO{} & \TODO{} \\
office\_s1\_06h (descr.)         & M1 ratio        & \TODO{} & \TODO{} & -- \\
\bottomrule
\end{tabular}
\end{adjustbox}
\end{table}

The camera ratios are 1.00 [0.94,\,1.06] and 1.07 [0.95,\,1.18]\src{T4}
\TODO{assign each to its sequence and state the interval type from the P9-D
results report (absent)}. At the 10\% margin\src{A2\,\S A2.2} one sequence is
non-inferior and the other inconclusive; H1 therefore formally fails
(subway\_s1\_24h M1 inconclusive\src{PCT}), and H3 is untested\src{PCT}. As
neither lower bound exceeds 1.0, this is reported as inconclusive, not as
degradation\src{A2\,\S A2.4}: we find no evidence of degradation.
\TODO{M2 and secondary metrics (results report).}

\begin{figure}[t]
\centering
\figplaceholder{\TODO{P9-E F1} (content per P9-E spec, absent)}{3.5cm}
\caption{\TODO{P9-E F1 caption.}}
\label{fig:f1}
\end{figure}

\subsection{Pedestrian results}
\begin{table}[t]
\centering
\caption{\TODO{P9-E T2 (pedestrians).} Medians over runs with pedestrian
samples\src{CR\,\S4}. Runs contributing: 36h ON 9, 36h LEV 6, 24h ON 12,
24h LEV 10 of 60\src{PCR\,\S1}. \TODO{Exact layout per P9-E spec.}}
\label{tab:ped}
\begin{adjustbox}{max width=\columnwidth}
\setlength{\tabcolsep}{3pt}\footnotesize
\begin{tabular}{@{}llcccc@{}}
\toprule
Cell & Metric & Graph & B0 & B1-causal & B1-RTS \\
\midrule
\multirow{2}{*}{subway\_s1\_36h} & pos.\ err.\ [m] & \TODO{} & \TODO{} & \TODO{} & \TODO{} \\
                                 & vel.\ err.\ [m/s] & \TODO{} & \TODO{} & \TODO{} & \TODO{} \\
\multirow{2}{*}{subway\_s1\_24h} & pos.\ err.\ [m] & \TODO{} & \TODO{} & \TODO{} & \TODO{} \\
                                 & vel.\ err.\ [m/s] & \TODO{} & \TODO{} & \TODO{} & \TODO{} \\
\midrule
\multicolumn{2}{@{}l}{H2a test statistic / $p$} & \multicolumn{4}{c}{\TODO{}} \\
\multicolumn{2}{@{}l}{H2b ratio [95\%]}         & \multicolumn{4}{c}{\TODO{}} \\
\bottomrule
\end{tabular}
\end{adjustbox}
\end{table}

H2 holds in both confirmatory cells\src{PCT}: velocity is better than
B1-causal (H2a) and position is no worse than B0 (H2b)\src{T4}.
\TODO{Numbers from the P9-D results report (absent).} The pre-registered
inclusion filter makes this a result on the best-aligned runs: the Sim(3)
RMSE $\leq\SI{0.5}{\metre}$ gate removes most subway runs (median per-map RMSE
\SIrange{1.27}{1.44}{\metre} on subway, \SIrange{0.17}{0.19}{\metre} on office),
leaving 6--12 of 60 runs per subway cell and arm\src{PCR\,\S1}.
\TODO{Post-hoc sensitivity analysis (pending\src{T4}).}

\begin{figure}[t]
\centering
\figplaceholder{\TODO{P9-E F2} (content per P9-E spec, absent)}{3.5cm}
\caption{\TODO{P9-E F2 caption.}}
\label{fig:f2}
\end{figure}

\begin{figure}[t]
\centering
\figplaceholder{\TODO{P9-E F3} (content per P9-E spec, absent)}{3.5cm}
\caption{\TODO{P9-E F3 caption.}}
\label{fig:f3}
\end{figure}

\subsection{Throughput}
\begin{table}[t]
\centering
\caption{Threaded-mode wall-clock on subway\_s1\_24h, one run per arm;
throughput only, not latency\src{PCR\,\S5}.}
\label{tab:timing}
\begin{adjustbox}{max width=\columnwidth}
\footnotesize
\begin{tabular}{@{}lcccc@{}}
\toprule
Arm & Wall [s] & Sensor [s] & Wall/sensor & Exit \\
\midrule
OFF & 117.5 & 738.0 & 0.159 & 0 \\
ON  &  99.8 & 738.0 & 0.135 & 0 \\
LEV & 121.6 & 738.0 & 0.165 & 0 \\
\bottomrule
\end{tabular}
\end{adjustbox}
\end{table}
Wall/sensor ratios of 0.135--0.165 correspond to about 6--7$\times$ real
time\src{PCR\,\S5}. The runner admits frames as fast as possible, so these are
offline throughput figures, not online latency; with one run per arm the
differences between arms are not interpretable\src{PCR\,\S5}. For comparison,
the deterministic campaign took a median of 1{,}317--1{,}700\,s wall
(0.43--0.56)\src{PCR\,\S5}.

\begin{figure}[t]
\centering
\figplaceholder{\TODO{P9-E F4} (content per P9-E spec, absent)}{3.5cm}
\caption{\TODO{P9-E F4 caption.}}
\label{fig:f4}
\end{figure}

% =====================================================================
\section{Discussion and Limitations}
\label{sec:limits}

\TODO{(handover) Discussion: what the analysis predicts and what the
experiments show; when pedestrians could help (near-zero static support).}

\paragraph*{Disclosures} Every disclosure below is carried from the
pre-registration and the post-check; \TODO{add every disclosure listed in the
handover (absent)}.
\begin{itemize}
  \item \emph{Reduced N and margins.} N and the margins were reduced from the
  Stage~0 plan (5\% margin, $N=260$ and 200) to $N=60$ and a 10\% / $+0.20$
  camera margin for compute reasons, before any P9-D data\src{A2\,\S A2.6}.
  \item \emph{Tuning data.} Tuning used apartment\_s1\_01h, apartment\_s1\_02h
  and office\_s1\_12h only; one evaluation-sequence pair was seen in
  P9-C\src{CR\,\S9}. \TODO{Name the pair (handover).}
  \item \emph{H2 on best-aligned runs.} The map-inclusion filter keeps 6--12 of
  60 subway runs per cell and arm\src{PCR\,\S1}; a post-hoc sensitivity
  analysis is pending\src{T4}.
  \item \emph{Camera.} H1 formally failed (inconclusive at the margin on
  subway\_s1\_24h); H3 was not tested\src{PCT}.
  \item \emph{Floor levels.} The levels module failed earlier gates (the
  synthetic gate, S6 spread; the shadow check; re-admission rule C) and was
  included by two overrides by Marco\src{CR\,\S1,\,\S9}; H4
  failed\src{PCT}. Not claimed.
  \item \emph{office\_s1\_06h.} Descriptive only: its mono-inertial
  initialisation fails with stock ORB-SLAM3 parameters\src{CR\,\S9}. Its OFF arm
  differed from Stage~0 (exit code 1 in 12/20 vs.\ 5/20 runs, Fisher two-sided
  $p=0.054$) with identical binaries, configs and command lines and only the
  settings-file text differing (layout sensitivity)\src{PCR\,\S4}. LEV
  coverage of 0.055 comes from a reset of the large first map at sensor time
  \SIrange{501.0}{502.6}{\second} in 11 of 20 LEV runs; no causal attribution is
  made\src{PCR\,\S4}.
  \item \emph{Determinism.} Deterministic mode is not stock ORB-SLAM3; three
  vanilla uninitialised-read contexts remain\src{CR\,\S9}.
  \item \emph{Terrain.} The subway sequences are multi-level: the camera
  changes height by about \SI{6}{\metre} and about 16--18\% of pedestrian
  samples are on non-level ground\src{CR\,\S9}.
  \item \emph{Analysis defects in descriptive columns.} Two descriptive columns
  of the frozen analysis script are wrong by construction (runs without an
  included map printed as 0; transitional share printed as 0 while 3.0--4.8\%
  of contact rows are flagged); neither reaches H2 or H4\src{PCR\,\S1--2}.
  \item \emph{Pre-run deviation.} The runner exit code of the five D1-P LEV runs
  on an evaluation sequence was read, slightly beyond the pre-registered
  list; no metric was read\src{PR\,\S7}.
  \item \emph{Simulation only.} Pedestrian accuracy on real data is not
  claimed\src{T4}; \TODO{InCrowd-VI / real-data plan (handover)}.
  \item \emph{Latency.} Not measured; only throughput\src{T4,\,PCR\,\S5}.
\end{itemize}

% =====================================================================
\section{Conclusion}
\label{sec:conclusion}
\TODO{(handover) Two to three sentences restating the supported claims at
their stated strength (Sec.~\ref{sec:intro}), the analysis, and the follow-up
on floor levels.}

% =====================================================================
\section*{Acknowledgment}
\TODO{(handover)}

% No references are cited yet: refs.bib is empty until the T1 report
% (verified sources) and the handover supply them.
\bibliographystyle{IEEEtran}
\bibliography{refs}

\end{document}
